\documentclass[11pt]{article}

\usepackage{color}
%\input{rgb}
%----------Packages----------
\usepackage{amsmath}
\usepackage{amssymb}
\usepackage{amsthm}
\usepackage{amsrefs}
\usepackage{dsfont}
\usepackage{mathrsfs}
\usepackage{stmaryrd}
\usepackage[all]{xy}
\usepackage[mathcal]{eucal}
\usepackage{verbatim}  %%includes comment environment
\usepackage{fullpage}  %%smaller margins
%----------Commands----------

%%penalizes orphans
\clubpenalty=9999
\widowpenalty=9999





%% bold math capitals
\newcommand{\bA}{\mathbf{A}}
\newcommand{\bB}{\mathbf{B}}
\newcommand{\bC}{\mathbf{C}}
\newcommand{\bD}{\mathbf{D}}
\newcommand{\bE}{\mathbf{E}}
\newcommand{\bF}{\mathbf{F}}
\newcommand{\bG}{\mathbf{G}}
\newcommand{\bH}{\mathbf{H}}
\newcommand{\bI}{\mathbf{I}}
\newcommand{\bJ}{\mathbf{J}}
\newcommand{\bK}{\mathbf{K}}
\newcommand{\bL}{\mathbf{L}}
\newcommand{\bM}{\mathbf{M}}
\newcommand{\bN}{\mathbf{N}}
\newcommand{\bO}{\mathbf{O}}
\newcommand{\bP}{\mathbf{P}}
\newcommand{\bQ}{\mathbf{Q}}
\newcommand{\bR}{\mathbf{R}}
\newcommand{\bS}{\mathbf{S}}
\newcommand{\bT}{\mathbf{T}}
\newcommand{\bU}{\mathbf{U}}
\newcommand{\bV}{\mathbf{V}}
\newcommand{\bW}{\mathbf{W}}
\newcommand{\bX}{\mathbf{X}}
\newcommand{\bY}{\mathbf{Y}}
\newcommand{\bZ}{\mathbf{Z}}

%% blackboard bold math capitals
\newcommand{\bbA}{\mathbb{A}}
\newcommand{\bbB}{\mathbb{B}}
\newcommand{\bbC}{\mathbb{C}}
\newcommand{\bbD}{\mathbb{D}}
\newcommand{\bbE}{\mathbb{E}}
\newcommand{\bbF}{\mathbb{F}}
\newcommand{\bbG}{\mathbb{G}}
\newcommand{\bbH}{\mathbb{H}}
\newcommand{\bbI}{\mathbb{I}}
\newcommand{\bbJ}{\mathbb{J}}
\newcommand{\bbK}{\mathbb{K}}
\newcommand{\bbL}{\mathbb{L}}
\newcommand{\bbM}{\mathbb{M}}
\newcommand{\bbN}{\mathbb{N}}
\newcommand{\bbO}{\mathbb{O}}
\newcommand{\bbP}{\mathbb{P}}
\newcommand{\bbQ}{\mathbb{Q}}
\newcommand{\bbR}{\mathbb{R}}
\newcommand{\bbS}{\mathbb{S}}
\newcommand{\bbT}{\mathbb{T}}
\newcommand{\bbU}{\mathbb{U}}
\newcommand{\bbV}{\mathbb{V}}
\newcommand{\bbW}{\mathbb{W}}
\newcommand{\bbX}{\mathbb{X}}
\newcommand{\bbY}{\mathbb{Y}}
\newcommand{\bbZ}{\mathbb{Z}}

%% script math capitals
\newcommand{\sA}{\mathscr{A}}
\newcommand{\sB}{\mathscr{B}}
\newcommand{\sC}{\mathscr{C}}
\newcommand{\sD}{\mathscr{D}}
\newcommand{\sE}{\mathscr{E}}
\newcommand{\sF}{\mathscr{F}}
\newcommand{\sG}{\mathscr{G}}
\newcommand{\sH}{\mathscr{H}}
\newcommand{\sI}{\mathscr{I}}
\newcommand{\sJ}{\mathscr{J}}
\newcommand{\sK}{\mathscr{K}}
\newcommand{\sL}{\mathscr{L}}
\newcommand{\sM}{\mathscr{M}}
\newcommand{\sN}{\mathscr{N}}
\newcommand{\sO}{\mathscr{O}}
\newcommand{\sP}{\mathscr{P}}
\newcommand{\sQ}{\mathscr{Q}}
\newcommand{\sR}{\mathscr{R}}
\newcommand{\sS}{\mathscr{S}}
\newcommand{\sT}{\mathscr{T}}
\newcommand{\sU}{\mathscr{U}}
\newcommand{\sV}{\mathscr{V}}
\newcommand{\sW}{\mathscr{W}}
\newcommand{\sX}{\mathscr{X}}
\newcommand{\sY}{\mathscr{Y}}
\newcommand{\sZ}{\mathscr{Z}}


\renewcommand{\phi}{\varphi}

\renewcommand{\emptyset}{\O}

\providecommand{\abs}[1]{\lvert #1 \rvert}
\providecommand{\norm}[1]{\lVert #1 \rVert}


\providecommand{\ar}{\rightarrow}
\providecommand{\arr}{\longrightarrow}

\renewcommand{\_}[1]{\underline{ #1 }}


\DeclareMathOperator{\ext}{ext}

\newcommand{\head}[1]{
	\begin{center}
		{\large #1}
		\vspace{.2 in}
	\end{center}
	
	\bigskip 
}

%----------Theorems----------

\newtheorem{theorem}{Theorem}[section]
\newtheorem{proposition}[theorem]{Proposition}
\newtheorem{lemma}[theorem]{Lemma}
\newtheorem{corollary}[theorem]{Corollary}


\newtheorem{axiom}{Axiom}


\theoremstyle{definition}
\newtheorem{definition}[theorem]{Definition}
\newtheorem{nondefinition}[theorem]{Non-Definition}
\newtheorem{exercise}[theorem]{Exercise}
\newtheorem{remark}[theorem]{Remark}
\newtheorem{remarks}[theorem]{Remarks}
\newtheorem{warning}[theorem]{Warning}


\numberwithin{equation}{subsection}


%----------Title-------------


\begin{document}

\head{MATH 161, Autumn 2025\\  SCRIPT 3: Introducing a Continuum }




This sheet introduces a continuum $C$ through a series of axioms.


\setcounter{section}{3}   

\begin{axiom} A continuum is a nonempty set $C$.  
\end{axiom}

We often refer to elements of $C$ as \emph{points}.


\begin{definition}  Let $X$ be a set.  An \emph{ordering} on the set $X$ is a subset $<$ of $X \times X$, with elements $(x, y) \in <$ written as $x < y$, satisfying the following properties:

\begin{itemize}
\item[(a)] {\it (Trichotomy) }

For all $x, y \in X$ exactly one of the following holds:
 $x < y,\  y < x\   \text{or } x=y.$ 
\item[(b)] {\it (Transitivity)} For all $x, y, z \in X$, if $x < y$ and $y < z$ then $x < z$.
\end{itemize}
\end{definition}


\begin{remarks}
\begin{enumerate}
\item[a)]  In mathematics ``or'' is understood to be inclusive unless stated otherwise. So in a) above, the word ``exactly'' is needed.
\item[b)] $x<y$ may also be written as $y>x.$
\item[c)] By $x \leq y$, we mean $x < y$ or $x = y$;  similarly for $x \geq y$.
\end{enumerate}
\end{remarks}


\begin{axiom}  A continuum $C$ has an ordering $<$.
\end{axiom}



\begin{definition}  If $A \subset C$ is a subset of $C$, then a point $a \in A$ is a \emph{first} point of $A$ if, for every element $x \in A$, either $a < x$ or $a = x$.  Similarly, a point $b \in A$ is called a \emph{last} point of $A$ if, for every $x \in A$, either $x < b$ or $x = b$.
\end{definition}

\begin{lemma}  If $A$ is a nonempty, finite subset of a continuum $C$, then $A$ has a first and last point.
\end{lemma}

\begin{theorem}  Suppose that $A$ is a set of $n$ distinct points in a continuum $C$, or, in other words, $A \subset C$ has cardinality $n$.  Then symbols $a_1, \dotsc, a_n$ may be assigned to each point of $A$ so that $a_1 < a_2 < \dotsm < a_n$, i.e. $a_i < a_{i + 1}$ for $1 \leq i \leq n - 1$.
\end{theorem}

\begin{definition}  If $x, y, z \in C$ and either (i) both $x < y$ and $y < z$ or (ii) both $z<y$ and $y<x,$ then we say that $y$ is \emph{between} $x$ and $z$.
\end{definition}

\begin{corollary}  Of three distinct points in a continuum, one must be between the other two.
\end{corollary}

\begin{axiom} A continuum $C$ has no first or last point.
\end{axiom}


In the next exercise we show that the integers and the rationals both have orderings that satisfy Axioms 1-3.

\begin{exercise}
\begin{enumerate}
\item[a)]  We define a relation $<$ on $\bbZ$ by $m<n$ if $n=m+c$ for some $c\in \bbN.$  Show that,
$\bbZ,$ with the ordering $<,$ satisfies Axioms 1-3. 

\item[b)] Show that, for any $p=\left[\frac{a}{b}\right]\in\bbQ,$ there is some $(a_1,b_1)\in p$ with $0<b_1.$

\item[c)] We define a relation $<_{\bbQ}$ on $\bbQ$ as follows. 
For $p,q\in\bbQ,$ let $(a_1,b_1)\in p$ be such that $0<b_1,$ and let $(a_2,b_2)\in q$ be such that $0<b_2.$  Then we define $p<_{\bbQ} q$ if $a_1b_2<a_2 b_1.$ 
Show that $<_{\bbQ}$ is a well-defined relation on $\bbQ.$
\item[d)] Show that $\bbQ,$ with the ordering $<_{\bbQ},$ satisfies Axioms 1-3.

\end{enumerate}
\end{exercise}

\begin{definition} If $a,b\in C$ and $a < b$, then the set of points between $a$ and $b$ is called a \emph{region}, denoted by $\_{ab}$.  
\end{definition}

\begin{warning}  One often sees the notation $(a, b)$ for regions.  We will reserve the notation $(a, b)$ for ordered pairs in a product $A \times B$.  These are very different things.  
\end{warning}

\begin{theorem} If $x$ is a point of a continuum $C$, then there exists a region $\_{ab}$ such that $x \in \_{ab}$.
\end{theorem}

We now come to one of the most important definitions of this course:

\begin{definition}
Let $A$ be a subset of a continuum $C$.  A point $p$ of $C$ is called a \emph{limit point} of $A$ if every region $R$ containing $p$ has nonempty intersection with $A \setminus \{p\}$.  Explicitly, this means:
\[
\text{for every region $R$ with $p \in R$, we have $R \cap (A \setminus \{p \}) \neq \emptyset$.}
\]
\end{definition}

Notice that we do not require that a limit point $p$ of $A$ be an element of $A$. We will use the notation $LP(A)$ to denote the set of limit points of $A.$ 

\begin{theorem} If $A \subset B$, then $LP(A)\subset LP(B).$
\end{theorem}



\begin{definition} If $\_{ab}$ is a region in a continuum $C$, then $C \setminus (\{a\} \cup \_{ab} \cup \{b\})$ is called the \emph{exterior} of $\_{ab}$ and is denoted by $\ext{\_{ab}}$.
\end{definition}


\begin{lemma} If $\_{ab}$ is a region in a continuum $C,$ then
$$\ext{\_{ab}}=\{x\in C\mid x<a\}\cup\{x\in C\mid b<x\}.$$
\end{lemma}



\begin{lemma}  No point in the exterior of a region is a limit point of that region.  No point of a region is a limit point of the exterior of that region.
\end{lemma}


\begin{theorem}  If two regions have a point $x$ in common, their intersection is a region containing $x$.
\end{theorem}

\begin{corollary}  If $n$ regions $R_1, \dotsc, R_n$ have a point $x$ in common, then their intersection $R_1 \cap \dotsm \cap R_n$ is a region containing $x$.
\end{corollary}

\begin{theorem}  Let $A, B$ be subsets of a continuum $C$.  Then  $LP(A \cup B)=LP(A)\cup LP(B).$
\end{theorem}



\begin{corollary}
Let $A_1, \dotsc, A_n$ be $n$ subsets of a continuum $C$.  Then  $p$ is a limit point of $(A_1 \cup \dotsm \cup A_n)$ if, and only if, $p$ is a limit point of at least one of the sets $A_k$.
\end{corollary}



\begin{theorem}  If $p$ and $q$ are distinct points of a continuum $C$, then there exist disjoint regions $R$ and~$S$ containing $p$ and $q$, respectively.
\end{theorem}

\begin{corollary}  A subset of a continuum $C$ consisting of one point has no limit points.
\end{corollary}

\begin{theorem} A finite subset $A$ of a continuum $C$ has no limit points.
\end{theorem}

\begin{corollary}  If $A$ is a finite subset of a continuum $C$ and $x \in A$, then there exists a region $R,$ containing $x,$ such that $A \cap R = \{ x \}$.
\end{corollary}


\begin{theorem}  If $p$ is a limit point of $A$ and $R$ is a region containing $p$, then the set $R \cap A$ is infinite.
\end{theorem}

\begin{exercise}  Find realizations of a continuum $(C, <)$.  That is, find concrete sets $C$ endowed with a relation $<$ satisfying all of the axioms (so far).  Are they the same?  What does ``the same'' mean here?
\end{exercise}
\bigskip

\begin{center}
{\em Additional Exercises}
\end{center} 

{\em In all exercises you are expected to prove your answer, unless explicitly stated otherwise.}

\begin{enumerate}

\item Show that for all $p,q\in \bbQ$ such that $p<_{\bbQ}q,$ there is some $r\in \bbQ$ with $p<_{\bbQ} r <_{\bbQ}q.$



\item By identifying $n\in \bbZ$ with $[\frac{n}{1}]\in \bbQ$ we can think of $\bbZ$ as a subset of $\bbQ.$ 
We shall write $n$ to mean $[\frac{n}{1}].$ Show that for all $p\in \bbQ,$ there is some 
$n\in \bbZ$ such that $p<_{\bbQ} n.$




\item


\begin{definition}
	Given two sets $A$ and $B$, we say that $A\subsetneq B$ if $A\subset B$ and $A\neq B$.
\end{definition}

For each set $X$ and subset $<_X \subset X\times X$, determine if $<_X$ is an ordering:
\begin{enumerate}
\item Let $X = \wp(\bbN)$ and $<_X = \{(A,B) \in \wp(\bbN)\times\wp(\bbN) | A \subset B\}$.
\item Let $X = \big\{ \{x \in \bbN | x \leq n\} \in \wp(\bbN) | n \in \bbN \big\}$ and $<_X = \{ (A, B)\in X\times X | A \subsetneq B\}$.
\item Let $X = \{ f \subset \bbN \times \bbN | f \text{ is a function}\}$ and $<_X = \{ (f,g) \in X\times X | f(n) < g(n) \text{ for all }n\in\bbN\}$.
\item Let $X = \{ f \subset \bbN \times \bbN | f \text{ is a function}\}$ and $<_X = \{ (f,g) \in X\times X | f(n) \leq g(n) \text{ for all }n\in\bbN \text{ and there exists } n\in\bbN \text{ such that } f(n) < g(n)\}$.
\end{enumerate}
(For the purposes of this exercise, ``$<$'' means the usual ordering on $\bbN$, i.e. if $m,n\in\bbN$ then $m < n$ if and only if $n = m+k$ for some $k\in\bbN$.)

\item
	Which of the following pairs of a set and an ordering satisfy Axioms 1,2, and 3 (and are thus examples of a ``continuum'')?
	\begin{enumerate}
		\item[i)] $C_1 = \{[n] \in \mathcal{P}(\bbN) \mid n \in \bbN\cup\{0\}\}$, where, for $[n], [m] \in C_1$, we say $[n] <_{C_1} [m]$ if $[n] \subsetneq [m]$;
		\item[ii)] $C_2 = \bbZ$, where $<_\bbZ$ is the usual ordering on $\bbZ$;
		\item[iii)] $C_3 = \bbZ\times \bbN$, where, for $(a,b), (x,y) \in \bbZ\times\bbN$ we say that $(a,b) <_3 (x,y)$ if $a < x$ or if $a=x$ and $b < y$.
	\end{enumerate}


\item
		Find, without proof, the exterior of each region.
		\begin{enumerate}
			\item Let $C_2$ and $<_\bbZ$ be as in Exercise 5.  Let $R = \underline{38}$.
			\item Let $C_3$ and $<_3$ be as in Exercise 5.  Let $R = \underline{(-2,5)(-2,10)}$.
			\item Let $C_3$ and $<_3$ be as in Exercise 5.  Let $R = \underline{(-2,5)(1,10)}$.
		\end{enumerate}



\item
		Find the limit points of each region.
		\begin{enumerate}
			\item Let $C_2$ and $<_\bbZ$ be as in Exercise 5.  Let $R = \underline{38}$.
			\item Let $C_3$ and $<_3$ be as in Exercise 5.  Let $R = \underline{(-2,5)(-2,10)}$.
			\item Let $C_3$ and $<_3$ be as in Exercise 5.  Let $R = \underline{(-2,5)(1,10)}$.
		\end{enumerate}


 

\item Let $A$ and $B$ be realizations of the continuum, with orderings $<_A, <_B,$ respectively. We say that $A$ and $B$ are {\em isomorphic} if there is a bijection $f:A\longrightarrow B$ such that 
$$a_1<_A a_2\Longrightarrow f(a_1)<_B f(a_2).$$
Show that $\bbZ$ and $\bbQ$ (with the orderings given in Scripts 1 and 2) are realizations of the continuum that are {\em not} isomorphic.




\item
Suppose that $R_1, R_2, \dots$ are regions in a continuum $C$ and that $x \in LP(\displaystyle \bigcup_{i=1}^\infty R_i)$.  Is it true that there exists $i \in \bbN$ such that $x \in LP(R_i)$?  Does this answer change if we add the additional hypothesis that there exists $k\in \bbN$ such that $x \in R_k$?

\item Show that any continuum $C$ satisfying Axioms 1-3 is infinite and write $C$ as a union of regions in $C.$



\item  Prove that if $A\subset\bbZ,$ $A $ has no limit points.



\item Let $i^2=-1$ and define $\bbZ [i]=\{a+bi\mid a,b\in\bbZ\}.$ Define an ordering on $\bbZ[i]$ by $a+bi<c+di$ if, and only if,
either $a<c,$ or $a=c$ and $b<d.$
\begin{enumerate}
\item Prove that $<$ is indeed an ordering. (So you must verify the properties in Definition 3.1. Note that you may assume the usual properties of $\bbZ$ - see Script 0 for full details.)
\item Prove that, with this ordering, $\bbZ[i]$ satisfies Axioms 1,2 and 3, so is a realization of the continuum as defined so far. 



\end{enumerate}




\end{enumerate}


\end{document}